\chapter[Reakcja--dyfuzja i fale biegnące]{Równania reakcji--dyfuzji, fale biegnące, dynamika długoczasowa i niestabilność Turinga}
\label{chap:pde27}

\section{Mechanizmy reakcji i dyfuzji: zakres rozdziału}
Rozważamy skalarne równanie
\begin{equation}\label{eq:pde27-main}
 u_t=D\Delta u+f(u),\qquad D>0,
\end{equation}
gdzie dyfuzja wygładza rozkład przestrzenny, a reakcja opisuje
lokalną dynamikę. Badamy zarówno torus
$\T_L^d=(\R/L\mathbb Z)^d$, jak i prostą $\R$.
Na torusie stosujemy periodyczne warunki brzegowe;
na prostej badamy fronty $u(t,x)=U(x-ct)$.
Jeżeli $f$ jest tylko lokalnie lipschitzowska, rozwiązanie
nie musi istnieć dla wszystkich czasów (por. rozdział~\ref{chap:pde25}).
W szczególności nie wolno wyprowadzać zachowania długoczasowego
z samego twierdzenia o lokalnym istnieniu.

W obliczeniach dla modelu Fishera--KPP użyjemy
$f(s)=rs(1-s)$, $r>0$. Rozwiązanie z zakresu $0\le u\le1$
pozostaje w tym przedziale, mimo że $f$ nie jest monotoniczna
na całej osi. Jest to zastosowanie zasady porównania,
nie monotoniczności funkcji $f$.

\section{Istnienie, ciągłość względem danych i zasada porównania}
\begin{twierdzenie}[Istnienie lokalne i alternatywa eksplozji]\label{thm:pde27-local}
Niech $L,D>0$, $d\ge1$, $f\in C^1(\R)$ oraz
$u_0\in C(\T_L^d)$. Istnieje dokładnie jedno maksymalne
rozwiązanie łagodne $u\in C([0,T_{\max});C(\T_L^d))$
równania~\eqref{eq:pde27-main} z $u(0)=u_0$.
Dla każdego $T<T_{\max}$ zachodzi wzór Duhamela
\[
 u(t)=S_D(t)u_0+\int_0^tS_D(t-s)f(u(s))\,\mathrm d\lambda_1(s),
 \qquad S_D(t)=e^{tD\Delta}.
\]
Jeżeli $T_{\max}<\infty$, to
$\|u(t)\|_{\infty}\to\infty$ przy $t\uparrow T_{\max}$.
Jeśli $u,v$ są rozwiązaniami o wartościach w $[-M,M]$
na $[0,T]$, to dla $K=\max_{|s|\le M}|f'(s)|$ mamy
\[
 \|u(t)-v(t)\|_{\infty}\le e^{Kt}\|u_0-v_0\|_{\infty},
 \qquad 0\le t\le T.
\]
\end{twierdzenie}
\begin{proof}
Na przestrzeni Banacha $X_T=C([0,T];C(\T_L^d))$
z normą $\|w\|_{X_T}=\sup_{0\le t\le T}\|w(t)\|_\infty$
rozważamy operator $\mathcal F$ prawej strony wzoru Duhamela,
w którym całka jest całką Bochnera w $C(\T_L^d)$.
Półgrupa cieplna jest silnie ciągła, dodatnia
i kontraktywna w normie maksimum; wynika to
z Lematu~\ref{lem:pde25-heat} przez skalowanie
w Stwierdzeniu~\ref{prop:pde-v44-torus-scaling}.
Dla $w\in X_T$ funkcja $s\mapsto f(w(s))$ jest ciągła
w normie maksimum, ponieważ zakres $w$ jest ograniczony.
Silna ciągłość półgrupy i jej kontraktywność dają
ciągłość całki względem $t$ (przez zbieżność
zmajoryzowaną na wspólnej części przedziałów całkowania
i oszacowanie całki na ich różnicy).
Zatem $\mathcal F$ przyjmuje wartości w $X_T$.

Ustalmy $M>\|u_0\|_\infty+1$ i oznaczmy
$F_M=\max_{|s|\le M}|f(s)|$, $K_M=\max_{|s|\le M}|f'(s)|$.
Dla $w,z$ z domkniętej kuli o promieniu $M$ mamy
\[
 \|\mathcal Fw\|_{X_T}\le\|u_0\|_\infty+TF_M,\qquad
 \|\mathcal Fw-\mathcal Fz\|_{X_T}
 \le TK_M\|w-z\|_{X_T}.
\]
Wybieramy $T>0$ tak, by $TF_M\le M-\|u_0\|_\infty$
i $TK_M<1$. Kula jest zupełna, więc twierdzenie
Banacha o punkcie stałym daje rozwiązanie.
Jednoznaczność w całej klasie rozwiązań łagodnych
wynika z oszacowania Gronwalla poniżej, zastosowanego
na dowolnym wspólnym zwartym przedziale istnienia.
Zgodne na przecięciach przedłużenia wyznaczają
jedno rozwiązanie na sumie wszystkich przedziałów
istnienia, czyli na $[0,T_{\max})$.
Wzór Duhamela można rozpocząć w dowolnej chwili
$t_0<T_{\max}$:
\[
 u(t_0+t)=S_D(t)u(t_0)+
 \int_0^t S_D(t-s)f(u(t_0+s))\,\mathrm d\lambda_1(s).
\]
Otrzymujemy go przez rozdzielenie pierwotnej całki
w chwili $t_0$ i użycie własności półgrupy.
Czas lokalnego istnienia można wybrać jednostajnie
dla danych o ograniczonej normie: dla
$\|u(t_0)\|_\infty\le B$ przyjmujemy $M=B+2$
i stosujemy powyższe warunki na $T$.
Nie zależy on od momentu rozpoczęcia ewolucji. Gdyby przy skończonym $T_{\max}$
norma rozwiązania była ograniczona wzdłuż pewnego ciągu
tendującego do $T_{\max}$, można by uruchamiać kontrakcję
od tych chwil na wspólnym przedziale o dodatniej długości,
sprzecznie z maksymalnością. To dowodzi alternatywy.
Dla różnicy $w=u-v$ wzór Duhamela i twierdzenie o wartości
średniej dają
$\|w(t)\|_\infty\le\|w(0)\|_\infty+
K\int_0^t\|w(s)\|_\infty\,\mathrm d\lambda_1(s)$.
Lemat Gronwalla kończy dowód stabilności.
\end{proof}


\begin{stwierdzenie}[Skalowanie półgrupy cieplnej na torusie]
\label{prop:pde-v44-torus-scaling}
Niech $L,D>0$ i $d\ge1$. Dla $v\in C(\T_L^d)$ połóżmy
$(Rv)(\theta)=v(L\theta/(2\pi))$, $\theta\in\T^d$.
Wówczas $R:C(\T_L^d)\to C(\T^d)$ jest liniową
izometrią na całą przestrzeń i
\begin{equation}\label{eq:pde-v44-torus-scale}
 R\,e^{tD\Delta_x}\,R^{-1}
 =e^{\,tD(2\pi/L)^2\Delta_\theta},\qquad t\ge0.
\end{equation}
Po podstawieniu $\tau=D(2\pi/L)^2t$ oraz
$w(\tau,\theta)=u(t,L\theta/(2\pi))$ równanie
$u_t=D\Delta_xu+f(u)$ przyjmuje postać
$w_\tau=\Delta_\theta w+\bigl[D(2\pi/L)^2\bigr]^{-1}f(w)$.
W konsekwencji konstrukcja lokalna i alternatywa eksplozji
z rozdziału~\ref{chap:pde25} stosują się do twierdzenia
\ref{thm:pde27-local} po tej zmianie zmiennych.
\end{stwierdzenie}
\begin{proof}
Skalowanie argumentu jest dobrze określone na klasach
modulo odpowiedni okres, a odwrotność ma postać
$(R^{-1}w)(x)=w(2\pi x/L)$. Obie normy maksimum
są równe, ponieważ zmiana zmiennych jest bijekcją torusów.
Z reguły łańcuchowej dla funkcji gładkiej
$R(\Delta_xv)=(2\pi/L)^2\Delta_\theta(Rv)$.
Dla wielomianów trygonometrycznych tożsamość półgrup
wynika także z działania na każdej częstotliwości
$k\in\mathbb Z^d$: obie strony mnożą współczynnik
Fouriera przez $\exp[-tD(2\pi/L)^2|k|^2]$.
Gęstość wielomianów trygonometrycznych w $C(\T^d)$
i kontraktywność półgrup rozszerzają równość na każdą
funkcję ciągłą. Zmiana czasu daje ostatni wzór.
Zmiana czasu we wzorze Duhamela daje tę samą
transformację także dla rozwiązań łagodnych.
Ponieważ zmodyfikowana nieliniowość nadal należy do
$C^1(\R)$, stosujemy Twierdzenie~\ref{thm:pde25-local}.
Ponadto Twierdzenie~\ref{thm:pde25-parabolic-regularity}
daje dla $0<a<T<T_{\max}$
\[
 u\in C([a,T];C^2(\T_L^d))\cap C^1([a,T];C(\T_L^d)),
\]
a równanie~\eqref{eq:pde27-main} jest spełnione
klasycznie dla dodatnich czasów.
\end{proof}


\begin{twierdzenie}[Zasada porównania i przedziały niezmiennicze]
\label{thm:pde27-comparison}
Niech $f\in C^1(\R)$, a $u,v$ będą rozwiązaniami z
poprzedniego twierdzenia, istniejącymi na $[0,T]$.
Jeżeli $u_0\le v_0$, to $u(t,\cdot)\le v(t,\cdot)$
dla każdego $0\le t\le T$. Jeśli $a\le b$,
$f(a)\ge0$ i $f(b)\le0$, to dane $a\le u_0\le b$
implikują $a\le u(t,\cdot)\le b$ na całym
przedziale istnienia; wówczas $T_{\max}=\infty$.
\end{twierdzenie}
\begin{proof}
Dzięki regularności dla dodatnich czasów, wykazanej
po Stwierdzeniu~\ref{prop:pde-v44-torus-scaling},
nie potrzebujemy gładkich danych początkowych.
Różnica $w=u-v$ spełnia
$w_t-D\Delta w=q(t,x)w$, gdzie
\[
 q(t,x)=\int_0^1 f'\bigl(v(t,x)+\theta(u(t,x)-v(t,x))\bigr)
 \,\mathrm d\lambda_1(\theta).
\]
Współczynnik $q$ jest ciągły i ograniczony na
$[0,T]\times\T_L^d$, gdyż $u,v$ są ciągłe na tym zbiorze.
Wybierzmy $K>\sup q$ i połóżmy $z=e^{-Kt}w$.
Wtedy $z_t-D\Delta z=(q-K)z$ dla $t>0$.
Gdyby $z$ przyjmowała wartość dodatnią, na zwartym
walcu $[0,T]\times\T_L^d$ osiągałaby dodatnie maksimum
w $(t_*,x_*)$ z $t_*>0$, ponieważ $z(0,\cdot)\le0$.
W tym punkcie $z_t\ge0$ (pochodna lewostronna,
jeśli $t_*=T$) i $\Delta z\le0$, natomiast
$(q-K)z<0$. Sprzeczność daje $z\le0$, czyli $u\le v$.

Ten sam argument działa, gdy $u$ jest podrozwiązaniem,
a $v$ nadrozwiązaniem, ciągłymi do chwili $0$
i klasycznymi dla $t>0$: równanie dla $z$
zastępujemy nierównością $z_t-D\Delta z\le(q-K)z$.
Stałe $a,b$ są odpowiednio pod- i nadrozwiązaniami,
gdyż $0\le f(a)$ i $0\ge f(b)$. Porównujemy je z $u$
na każdym zwartym przedziale jego istnienia.
Przedział $[a,b]$ jest niezmienniczy, a ograniczenie normy
wyklucza skończony czas maksymalny.
\end{proof}

\begin{kontrprzyklad}[Porównanie nie wymaga $f'\ge0$]
Dla $f(s)=s(1-s)$ pochodna $f'(s)=1-2s$
jest ujemna na $(1/2,\infty)$.
Mimo to zasada porównania działa, ponieważ różnica
$f(u)-f(v)$ ma postać $a(t,x)(u-v)$ z ograniczonym
współczynnikiem na każdym ograniczonym przedziale czasu.
Dla układów reakcji--dyfuzji takie porównanie
składowa po składowej na ogół nie zachodzi bez
warunku kooperatywności.
\end{kontrprzyklad}

\section{Równanie Fishera--KPP: zbieżność do stanu nasycenia na torusie}
\begin{twierdzenie}[Zbieżność do stanu nasycenia]\label{thm:pde27-kpp-torus}
Niech $u_0\in C(\T_L^d)$, $0\le u_0\le1$ i
$u_0\not\equiv0$. Dla $D,r>0$ jedyne rozwiązanie
zagadnienia
$u_t=D\Delta u+ru(1-u)$ z danymi $u_0$
istnieje globalnie i
\[
 \lim_{t\to\infty}\|u(t,\cdot)-1\|_{C(\T_L^d)}=0.
\]
Jeśli $u_0\equiv0$, to $u(t,\cdot)\equiv0$.
\end{twierdzenie}
\begin{proof}
Porównanie ze stałymi rozwiązaniami $0$ i $1$
daje $0\le u\le1$ i istnienie globalne.
Dla $t>0$ półgrupa $S_D(t)$ ma dodatnie jądro
periodyczne względem miary Lebesgue'a:
\[
 K_{D,t}^L(z)=\sum_{k\in\mathbb Z^d}
 (4\pi Dt)^{-d/2}\exp\left(-\frac{|z+Lk|^2}{4Dt}\right).
\]
Wzór wynika z Lematu~\ref{lem:pde25-heat}
przez zmianę skali, łącznie z jakobianem miary.
Szereg jest jednostajnie zbieżny dla ustalonego $t>0$
na kostce fundamentalnej i wszystkie składniki są
dodatnie. Ponieważ $u_0$ jest nieujemna, ciągła
i niezerowa, jest dodatnia na zbiorze dodatniej miary.
Stąd $S_D(t)u_0>0$ punktowo.
Wzór Duhamela z nieujemnym źródłem $ru(1-u)$ daje
$u(t,\cdot)\ge S_D(t)u_0$, więc dla każdego $t_0>0$
funkcja $u(t_0,\cdot)$ jest ściśle dodatnia. Zwarta przestrzeń
$\T_L^d$ zapewnia minimum $m_0>0$.
Niech $y'=ry(1-y)$ i $z'=rz(1-z)$,
$y(t_0)=m_0$, $z(t_0)=\max u(t_0,\cdot)\le1$.
Porównanie daje $y(t)\le u(t,x)\le z(t)$.
Z jawnego rozwiązania równania logistycznego
\[
 y(t)=\frac{1}{1+(m_0^{-1}-1)e^{-r(t-t_0)}}
\]
wynika $y(t)\uparrow1$ (dla $m_0=1$ funkcja
jest stale równa $1$), podczas gdy $z(t)\le1$.
Wzór na $y$ sprawdzamy przez różniczkowanie
i podstawienie warunku początkowego; analogiczny
wzór obowiązuje dla $z$. Porównanie jest dopuszczalne
na każdym skończonym przedziale po chwili $t_0$,
gdyż oba rozwiązania skalarne są globalne.
Mamy $\|1-u(t)\|_\infty\le1-y(t)\to0$.
Dla $u_0\equiv0$ teza wynika z jednoznaczności.
\end{proof}

\begin{uwaga}[Dlaczego torus, a nie cała prosta?]
Na $\R$ rozwiązanie wychodzące z danych o zwartym nośniku
nie zbiega jednostajnie do $1$ na całej prostej
przy $t\to\infty$: dla każdego skończonego $t>0$
rozwiązanie dąży do zera przy $|x|\to\infty$, co
wynika z oszacowania gaussowskiego w dowodzie
Twierdzenia~\ref{thm:pde27-spreading-upper}
(dla $0\le u_0\le1$). W szczególności
$\|1-u(t)\|_\infty=1$. Właściwą asymptotyką jest
\emph{propagacja frontu}, z którą zapoznamy się dalej.
\end{uwaga}


\begin{wniosek}[Ilościowa zbieżność rozwiązania Fishera--KPP na torusie]
\label{cor:pde27-kpp-rate}
Przy hipotezach twierdzenia~\ref{thm:pde27-kpp-torus}
niech $t_0>0$ i
$m_{t_0}=\min_{x\in\T_L^d}(S_D(t_0)u_0)(x)>0$.
Wówczas dla wszystkich $t\ge t_0$
\begin{equation}\label{eq:pde27-kpp-rate}
 \|1-u(t,\cdot)\|_\infty\le
 \frac{1-m_{t_0}}{m_{t_0}}e^{-r(t-t_0)}.
\end{equation}
Ponadto jeśli $K_{D,t_0}^L$ oznacza dodatnie
periodyczne jądro cieplne na $\T_L^d$, to
\[
 m_{t_0}\ge\left(\min_{z\in\T_L^d}K_{D,t_0}^L(z)\right)
 \int_{\T_L^d}u_0(x)\,\mathrm d\lambda_d(x)>0.
\]
\end{wniosek}
\begin{proof}
Nieujemne źródło $ru(1-u)$ i wzór Duhamela dają
$u(t_0)\ge S_D(t_0)u_0$, stąd
$\min u(t_0)\ge m_{t_0}>0$.
Rozwiązanie skalarnego równania
$y'=ry(1-y)$ z $y(t_0)=m_{t_0}$ wynosi
$y(t)=(1+(m_{t_0}^{-1}-1)e^{-r(t-t_0)})^{-1}$.
Porównanie ze stałym w przestrzeni rozwiązaniem $y$
i nadrozwiązaniem $1$ daje
$0\le1-u(t,x)\le1-y(t)$.
Ponieważ $1-(1+a)^{-1}=a/(1+a)\le a$
dla $a\ge0$, otrzymujemy~\eqref{eq:pde27-kpp-rate}.
Ostatnia nierówność wynika bezpośrednio z dodatniości
i ciągłości jądra na zwartym torusie.
\end{proof}


\section{Fale biegnące: od PDE do układu równań zwyczajnych}
Niech $d=1$ i rozważmy równanie
$u_t=Du_{xx}+ru(1-u)$ na $\R$.
Dla postaci fali biegnącej $u(t,x)=U(z)$, $z=x-ct$,
mamy $u_t=-cU'$, $u_x=U'$ oraz $u_{xx}=U''$,
a zatem otrzymujemy równanie profilu
\begin{equation}\label{eq:pde27-wave}
 DU''+cU'+rU(1-U)=0.
\end{equation}
Szukamy funkcji $U\in C^2(\R)$, dla której
$U(-\infty)=1$, $U(+\infty)=0$ i $U'<0$.
Prędkość $c$ jest prędkością przemieszczania się
profilu w prawo. Sama zamiana zmiennych
nie dowodzi jeszcze istnienia takiego profilu.

\begin{lemat}[Obszar niezmienniczy dla układu profilowego]
\label{lem:pde27-wedge}
Niech $D,r>0$, $c\ge2\sqrt{Dr}$ oraz
\[
 k=\frac{c-\sqrt{c^2-4Dr}}{2D}>0.
\]
Wtedy obszar
$\mathcal W=\{(U,q):0<U<1,\ 0<q<kU\}$
jest dodatnio niezmienniczy dla układu
\begin{equation}\label{eq:pde27-phase}
 U'=-q,\qquad q'=-\frac cD q+\frac rD U(1-U).
\end{equation}
Oznacza to, że trajektorie rozpoczynające się w $\mathcal W$
pozostają w nim dla dodatnich czasów istnienia.
Ponadto rozwiązanie wychodzące z punktu
$U(0)=1-\delta$, $q(0)=0$, gdzie $0<\delta<1/2$,
zachowuje $0<U<1$, $0<q<kU$ dla $z>0$ i
spełnia $U(z)\to0$, $q(z)\to0$ przy $z\to\infty$.
\end{lemat}
\begin{proof}
Na brzegu $q=0$, gdy $0<U<1$, mamy
$q'=rU(1-U)/D>0$. Na brzegu $q=kU$ otrzymujemy
\[
 (q-kU)'=\bigl(k^2-\tfrac cDk+\tfrac rD(1-U)\bigr)U
 =-\tfrac rD U^2<0,
\]
ponieważ $k^2-(c/D)k+r/D=0$.
Uzasadnijmy brak pierwszego wyjścia z obszaru.
Dopóki trajektoria pozostaje w $\mathcal W$, mamy
$U'<0$ oraz $U'+kU>0$. Stąd
$U(0)e^{-kz}\le U(z)\le U(0)<1$.
Jeżeli pierwszy kontakt z brzegiem nastąpiłby
w skończonej chwili $z_*>0$, to powyższe oszacowanie
wykluczałoby brzegi $U=0$ i $U=1$.
Przy pierwszym kontakcie z $q=0$ od strony $q>0$
musiałoby zachodzić $q'(z_*)\le0$, wbrew
$q'(z_*)=rU(z_*)(1-U(z_*))/D>0$.
Przy pierwszym kontakcie z $q=kU$ od strony $q-kU<0$
musiałoby zachodzić $(q-kU)'(z_*)\ge0$, wbrew
obliczonej ujemnej wartości.
Trajektoria pozostaje zatem w $\mathcal W$.
Pole wektorowe jest wielomianowe, więc jest lokalnie
lipschitzowskie; ograniczoność trajektorii w zwartym
domknięciu $\mathcal W$ wyklucza skończony maksymalny
czas istnienia (lokalny czas istnienia można wybrać
jednostajnie dla danych z tego zwartego zbioru).
Dla danych $U(0)=1-\delta$, $q(0)=0$ mamy
$q'(0)>0$ i $q(0)-kU(0)<0$. Ciągłość daje
$0<U<1$ i $0<q<kU$ dla dostatecznie małych $z>0$,
po czym stosujemy poprzedni argument.
Funkcja $U$ jest malejąca, więc ma granicę
$\ell\in[0,1-\delta)$. Ponieważ $q=-U'\ge0$,
$\int_0^\infty q(z)\,\mathrm d\lambda_1(z)
=U(0)-\ell<\infty$, przez twierdzenie Newtona--Leibniza
na przedziałach skończonych i przejście do granicy.
Pochodna $q'$ jest ograniczona w $\mathcal W$,
więc $q$ jest jednostajnie ciągła.
Gdyby $q$ nie dążyła do zera, istniałby ciąg
$z_j\to\infty$ z $q(z_j)\ge\varepsilon>0$.
Jednostajna ciągłość dawałaby wspólny promień $\rho>0$
taki, że $q\ge\varepsilon/2$ na przedziałach
$[z_j-\rho,z_j+\rho]$. Wybierając rozłączny
podciąg tych przedziałów, otrzymalibyśmy nieskończoną
całkę z $q$. Zatem $q(z)\to0$. Gdyby $\ell>0$, drugie
równanie~\eqref{eq:pde27-phase} prowadziłoby do
$q'(z)\to r\ell(1-\ell)/D>0$, co jest sprzeczne
z $q(z)\to0$. W konsekwencji $\ell=0$.
\end{proof}

\begin{twierdzenie}[Prędkości fal biegnących równania Fishera--KPP]
\label{thm:pde27-wave}
Niech $D,r>0$. Dla każdej $c\ge2\sqrt{Dr}$
istnieje $U\in C^2(\R)$ spełniająca
\eqref{eq:pde27-wave}, $0<U<1$, $U'<0$ oraz
$U(-\infty)=1$, $U(+\infty)=0$.
Natomiast dla żadnej $c<2\sqrt{Dr}$
taki profil nie istnieje. W szczególności
minimalna prędkość monotonicznej fali wynosi
$c_*=2\sqrt{Dr}$.
\end{twierdzenie}
\begin{proof}
\textbf{Krok 1: skonstruowanie pełnej trajektorii.}
Ustalmy $c\ge2\sqrt{Dr}$ i $k$ z
lematu~\ref{lem:pde27-wedge}. Dla $\delta\downarrow0$
rozważmy rozwiązania $(U_\delta,q_\delta)$
rozpoczynające się w $(1-\delta,0)$.
Każda trajektoria osiąga poziom $U=1/2$ dokładnie raz
na mocy ciągłości, ścisłego malenia i granicy zero;
niech $t_\delta>0$ oznacza odpowiedni czas.
Pokażemy, że $t_\delta\to\infty$.
Kładziemy $v=1-U$ i korzystamy z równań
$v'=q$, $q'=-(c/D)q+(r/D)(1-v)v\le(r/D)v$.
Dla $v,q\ge0$ dostajemy
$(v+q)'\le C(v+q)$ z $C=\max\{1,r/D\}$,
a zatem $v(t)+q(t)\le\delta e^{Ct}$.
Gdy $U(t_\delta)=1/2$, mamy $v(t_\delta)=1/2$,
więc $t_\delta\ge C^{-1}\log(1/(2\delta))\to\infty$.

Przesuwamy trajektorie: $\widetilde U_\delta(z)
=U_\delta(t_\delta+z)$, analogicznie dla $q$.
Są one określone na $(-t_\delta,\infty)$,
mają wartości w zwartym domknięciu $\mathcal W$
i jednostajnie ograniczone pochodne.
Oznaczmy przez $F$ pole wektorowe układu
~\eqref{eq:pde27-phase}. Jest ono klasy $C^\infty$,
więc na zwartym domknięciu $\mathcal W$ zarówno $F$,
jak i $DF$ są ograniczone. Dla przesuniętych trajektorii
$Y_\delta=(\widetilde U_\delta,\widetilde q_\delta)$
mamy $Y_\delta'=F(Y_\delta)$ i
$Y_\delta''=DF(Y_\delta)F(Y_\delta)$.
Zatem na każdym $[-m,m]$ dla dostatecznie małych $\delta$
funkcje i ich pierwsze pochodne są jednostajnie
ograniczone i jednakowo ciągłe.
Wybierając najpierw ciąg $\delta_j\downarrow0$,
stosujemy twierdzenie Arzeli--Ascolego do par
$(Y_{\delta_j},Y_{\delta_j}')$ na $[-m,m]$
i wydzielamy podciąg diagonalny po $m\in\mathbb N$.
Jednostajna zbieżność funkcji i pochodnych
oraz twierdzenie Newtona--Leibniza dają zbieżność
w $C^1_{\mathrm{loc}}(\R)$. Przechodząc do granicy
w $Y_\delta'=F(Y_\delta)$, otrzymujemy pełną
trajektorię $(U,q)$ układu~\eqref{eq:pde27-phase}
o wartościach w $\overline{\mathcal W}$, z $U(0)=1/2$. Ponieważ $q\ge0$ i $U'=-q$,
$U$ jest nierosnąca i ma granice $\ell_\pm$
przy $z\to\pm\infty$. Analogicznie jak w lemacie,
całkowalność $q$ na obu półprostych i jego
jednostajna ciągłość implikują $q(z)\to0$
na obu krańcach. Drugie równanie układu daje
$q'(z)\to (r/D)\ell_\pm(1-\ell_\pm)$
przy $z\to\pm\infty$. Gdyby któraś z tych granic
była dodatnia, to $q'$ byłaby ograniczona od dołu
przez dodatnią stałą na odpowiedniej półprostej.
Całkowanie nierówności na dowolnie długim odcinku
przeczyłoby ograniczoności $q$. Stąd
$\ell_\pm(1-\ell_\pm)=0$.
Z $U(0)=1/2$ i monotoniczności wynika
$\ell_-=1$, $\ell_+=0$.
Gdyby $U(z_0)=0$ lub $U(z_0)=1$, to $U$
osiągałaby odpowiednio minimum lub maksimum
w punkcie wewnętrznym prostej, a zatem $q(z_0)=0$.
Oba punkty $(0,0)$ i $(1,0)$ są równowagami
układu. Jednoznaczność rozwiązania zagadnienia
początkowego w obu kierunkach czasu wymuszałaby
trajektorię stałą, wbrew $U(0)=1/2$.
Zatem $0<U<1$ na całej prostej.
Gdyby $q(z_0)=0$ przy $0<U(z_0)<1$,
to $q'(z_0)>0$, co jest niemożliwe przy $q\ge0$
na całej prostej. Stąd $q>0$ i $U'<0$.

\textbf{Krok 2: wykluczenie prędkości mniejszych od krytycznej.}
Jeżeli $c\le0$, po podstawieniu $q=-U'>0$
do~\eqref{eq:pde27-phase} otrzymujemy $q'>0$,
a to przeczy jednocześnie całkowalności
$\int_\R q(z)\,\mathrm d\lambda_1(z)=1$
(wynika to z granic $U$ i $q=-U'$):
dla dowolnego $a\in\R$ mamy wówczas
$q(z)\ge q(a)>0$ dla $z\ge a$, co przeczy
skończoności całki.
Niech teraz $0<c<2\sqrt{Dr}$. Dla $z$ dostatecznie dużych
$U(z)$ jest tak małe, że
$\mu(z)=\frac rD(1-U(z))-\frac{c^2}{4D^2}\ge\kappa^2>0$.
Po podstawieniu $y(z)=e^{cz/(2D)}U(z)>0$
równanie~\eqref{eq:pde27-wave} przyjmuje postać
$y''+\mu(z)y=0$. Istotnie, pisząc
$U=e^{-cz/(2D)}y$, dostajemy
\[
 DU''+cU'=e^{-cz/(2D)}
 \left(Dy''-\frac{c^2}{4D}y\right),
\]
a składnik reakcji po podzieleniu przez
$De^{-cz/(2D)}$ ma postać $(r/D)(1-U)y$.
Możemy wybrać na przykład
$\kappa^2=\frac12(r/D-c^2/(4D^2))>0$,
a następnie $a$ tak duże, by
$(r/D)U(z)\le\kappa^2$ dla $z\ge a$.
Na przedziale $[a,a+\pi/\kappa]$, gdzie $a$ jest
dostatecznie duże, mnożymy równanie przez
$\sin(\kappa(z-a))$ i całkujemy względem $z$.
Po dwukrotnym całkowaniu przez części dostajemy
\[
 \int_a^{a+\pi/\kappa}(y''+\kappa^2y)
 \sin(\kappa(z-a))\,\mathrm d\lambda_1(z)
 =\kappa\bigl(y(a)+y(a+\pi/\kappa)\bigr)>0.
\]
Tymczasem $y''+\kappa^2y=-(\mu-\kappa^2)y\le0$,
a sinus jest nieujemny na całym przedziale.
Sprzeczność kończy dowód.
\end{proof}

\begin{przyklad}[Jawna fala o prędkości $5/\sqrt6$]
\label{ex:pde27-explicit}
Dla $D=r=1$ i $c=5/\sqrt6$ profil
\[
 U(z)=\frac{1}{(1+e^{z/\sqrt6})^2}
\]
rozwiązuje~\eqref{eq:pde27-wave} i spełnia
$U(-\infty)=1$, $U(+\infty)=0$, $U'<0$.
Podstawiając $w=e^{z/\sqrt6}$ mamy
$U'=-2w/(\sqrt6(1+w)^3)$ oraz
$U''=-w/(3(1+w)^3)+w^2/(1+w)^4$;
po sprowadzeniu do wspólnego mianownika $(1+w)^4$
licznik wyrażenia $U''+(5/\sqrt6)U'+U(1-U)$
wynosi
$-\frac13 w(1+w)+w^2-\frac53w(1+w)+2w+w^2=0$.
Znak $U'$ i granice wynikają bezpośrednio ze wzoru.
Dla ogólnych $D,r>0$ piszemy
$U(z)=V(\sqrt{r/D}\,z)$, gdzie $V$ jest powyższym
profilem dla $D=r=1$. Po podstawieniu i podzieleniu
przez $r$ równanie profilu ma postać
$V''+(c/\sqrt{Dr})V'+V(1-V)=0$.
Wybór $c=(5/\sqrt6)\sqrt{Dr}$ daje zatem
rozwiązanie
\[
 U(z)=\left(1+\exp\left(\sqrt{\frac r{6D}}\,z\right)\right)^{-2}.
\]
\end{przyklad}

\begin{center}
\begin{tikzpicture}[>=Latex]
 \begin{scope}[x=4.8cm,y=4.0cm]
  \fill[blue!9] (0,0)--(1,0)--(1,.8164966)--cycle;
  \draw[->] (-.04,0)--(1.13,0) node[right] {$U$};
  \draw[->] (0,-.035)--(0,1.02) node[above] {$q$};
  \draw[blue!70!black,thick] (0,0)--(1,.8164966)
    node[above left] {$q=kU$};
  \draw[dashed] (1,0)--(1,.8164966);
  \draw[->,thick] (.5,0)--(.5,.12);
  \draw[->,thick] (.68,.5552177)--(.58,.42);
  \draw[->,thick] (1,.45)--(.88,.40);
  \draw[red!75!black,very thick,domain=0:1,samples=101]
    plot (\x,{2/sqrt(6)*\x*(1-sqrt(\x))});
  \draw[red!75!black,->,thick] (.56,{2/sqrt(6)*.56*(1-sqrt(.56))})
    --(.46,{2/sqrt(6)*.46*(1-sqrt(.46))});
  \fill (0,0) circle (1.3pt) node[below left] {$0$};
  \fill (1,0) circle (1.3pt) node[below] {$1$};
  \node at (.52,.96) {Obszar fazowy};
 \end{scope}
 \begin{scope}[xshift=6.9cm,x=.43cm,y=3.0cm]
  \draw[->] (-6.5,0)--(6.7,0) node[right] {$z$};
  \draw[->] (0,-.05)--(0,1.19) node[above] {$U$};
  \draw[dashed,gray] (-6.3,1)--(6.3,1);
  \node[left] at (0,1) {$1$};
  \draw[red!75!black,very thick,domain=-6.3:6.3,samples=101]
    plot (\x,{1/(1+exp(\x/sqrt(6)))^2});
  \node[below] at (-6,0) {$-6$};
  \node[below] at (6,0) {$6$};
  \node at (0,1.36) {Profil fali};
 \end{scope}
\end{tikzpicture}

\small Rysunek: obszar niezmienniczy $\mathcal W$ i jawna
trajektoria $q=(2/\sqrt6)U(1-\sqrt U)$ (czerwona)
dla $D=r=1$, $c=5/\sqrt6$, $k=2/\sqrt6$.
Po prawej odpowiadający jej profil $U(z)$.
\end{center}

\section{Granice propagacji na prostej}
\begin{twierdzenie}[Górne oszacowanie prędkości ekspansji]
\label{thm:pde27-spreading-upper}
Niech $D,r>0$, a $u_0\in C_c(\R)$ spełnia
$0\le u_0\le1$. W przestrzeni $\operatorname{BUC}(\R)$ rzeczywistych
ograniczonych funkcji jednostajnie ciągłych istnieje
dokładnie jedno globalne rozwiązanie łagodne
$u\in C([0,\infty);\operatorname{BUC}(\R))$
równania $u_t=Du_{xx}+ru(1-u)$ z danymi $u_0$.
Spełnia ono $0\le u\le1$.
Wówczas dla każdej $c>2\sqrt{Dr}$ istnieją
stałe $C,\gamma>0$, zależne od $c,u_0,D,r$, takie, że
\[
 \sup_{|x|\ge ct}u(t,x)\le C e^{-\gamma t}
 \quad\text{dla wszystkich }t\ge0.
\]
\end{twierdzenie}
\begin{proof}
\textbf{Istnienie i porównanie na prostej.}
Przestrzeń $\operatorname{BUC}(\R)$ jest domkniętą
podprzestrzenią przestrzeni funkcji ograniczonych
z normą supremum, a więc jest przestrzenią Banacha.
Dla $t>0$ oznaczmy przez $S_D(t)$ splot z jądrem
$G_{Dt}(x)=(4\pi Dt)^{-1/2}\exp(-x^2/(4Dt))$,
a $S_D(0)=I$. Jądro jest nieujemne i ma całkę $1$.
Z tożsamości splotowej jąder Gaussa wynika własność
półgrupy; całkowanie daje kontraktywność, dodatniość
i zachowanie stałych. Splot zachowuje jednostajną
ciągłość, ponieważ
$\|S_D(t)v(\cdot+h)-S_D(t)v\|_\infty
\le\|v(\cdot+h)-v\|_\infty$.
Silna ciągłość przy $t=0$ wynika z rozdzielenia
całki z $G_{Dt}(y)(v(x-y)-v(x))$ na $|y|<\eta$
i $|y|\ge\eta$: na pierwszym zbiorze stosujemy
jednostajną ciągłość $v$, a na drugim zanik masy
ogona Gaussa przy $t\downarrow0$.
Silna ciągłość w pozostałych chwilach wynika
z własności półgrupy i kontraktywności.

Niech $\pi(s)=\min\{1,\max\{0,s\}\}$ oraz
$g(s)=r\pi(s)(1-\pi(s))$. Funkcja $g$ jest globalnie
lipschitzowska ze stałą $r$. Dla każdego $T>0$
na $C([0,T];\operatorname{BUC}(\R))$ rozważamy
operator
\[
 (\mathcal Pv)(t)=e^{-rt}S_D(t)u_0+
 \int_0^t e^{-r(t-s)}S_D(t-s)
       (g(v(s))+rv(s))\,\mathrm d\lambda_1(s).
\]
Ciągłość całki i kontrakcję uzasadniamy tak samo
jak w dowodzie Twierdzenia~\ref{thm:pde27-local}.
Dla $2rT<1$ operator jest kontrakcją na całej
przestrzeni, ponieważ $g(s)+rs$ ma stałą
Lipschitza co najwyżej $2r$. Funkcja $g(s)+rs$
jest niemalejąca z nierówności Lipschitza.
Stąd $\mathcal P$ zachowuje porządek punktowy.
Dla funkcji stale równych $0$ i $1$ mamy
$\mathcal P0\ge0$ oraz $\mathcal P1\le1$,
ponieważ $0\le u_0\le1$, $g(0)=g(1)=0$
i $S_D(t)$ zachowuje stałe.
Iteracje Picarda rozpoczęte od zera pozostają
w przedziale $[0,1]$; ich granica również.
Powtarzając konstrukcję na kolejnych przedziałach
tej samej długości, otrzymujemy rozwiązanie globalne.

W przedziale $[0,1]$ funkcja $g$ jest równa
$f(s)=rs(1-s)$. Wzór z $\mathcal P$ jest równoważny
zwykłemu wzorowi Duhamela z półgrupą $S_D$.
Istotnie, dla dowolnej rzeczywistej stałej $a$
mamy
\[
 e^{at}S_D(t)=S_D(t)+a\int_0^t
 S_D(t-s)e^{as}S_D(s)\,\mathrm d\lambda_1(s).
\]
Podstawienie tej tożsamości do wzoru Duhamela
i zastosowanie twierdzenia Fubiniego do całek
Bochnera (na skończonych przedziałach całkowane
funkcje są ograniczone w normie) pokazują,
że przejściu od $S_D$ do $e^{at}S_D(t)$
towarzyszy zastąpienie $f(u)$ przez $f(u)-au$.
Odwrotna transformacja odpowiada stałej $-a$.
Jednoznaczność w całej klasie rozwiązań łagodnych
dla $f$ wynika z oszacowania Gronwalla na każdym
skończonym przedziale, na którym oba rozwiązania
mają ograniczone normy supremum.

\textbf{Oszacowanie gaussowskie.}
Stosujemy powyższą transformację z $a=r$.
Ponieważ $f(u)-ru=-ru^2$, otrzymujemy
\[
 u(t)=e^{rt}S_D(t)u_0-
 r\int_0^t e^{r(t-s)}S_D(t-s)u(s)^2\,
                       \mathrm d\lambda_1(s).
\]
Dodatniość półgrupy i wcześniej wykazane
$0\le u\le1$ dają
\[
 0\le u(t,x)\le e^{rt}\int_{\R}
 \frac{e^{-(x-y)^2/(4Dt)}}{\sqrt{4\pi Dt}}u_0(y)\,\mathrm d\lambda_1(y),
 \qquad t>0.
\]
Wybierzmy $R>0$ tak, aby $\supp u_0\subset[-R,R]$,
oraz liczbę $c_0$ z $2\sqrt{Dr}<c_0<c$.
Niech $t_0=\max\{1,R/(c-c_0)\}$. Dla $t\ge t_0$
i $|x|\ge ct$ zachodzi $|x-y|\ge c_0t$
dla każdego $y\in\supp u_0$. Stąd
\[
 u(t,x)\le\frac{\|u_0\|_{L^1(\R)}}{\sqrt{4\pi Dt}}
 \exp\!\left(-\left(\frac{c_0^2}{4D}-r\right)t\right).
\]
Liczba $\gamma=c_0^2/(4D)-r$ jest dodatnia,
a czynnik $t^{-1/2}\le1$ dla $t\ge t_0$.
Na odcinku $0\le t<t_0$ korzystamy z $u\le1$.
Wystarczy więc przyjąć
$C=\max\{e^{\gamma t_0},\|u_0\|_{L^1}/\sqrt{4\pi D}\}$,
co daje żądane oszacowanie dla wszystkich $t\ge0$.
\end{proof}

\begin{uwaga}[Co udowodniono, a czego nie]
Twierdzenie~\ref{thm:pde27-spreading-upper} daje wyłącznie
górne oszacowanie szybkości rozprzestrzeniania się.
Dopasowana dolna granica prędkości i pełne twierdzenie
o asymptotyce rozwiązania na całej prostej wymagają
osobnej konstrukcji podrozwiązań. Nie wynikają
same z istnienia fal ani z powyższej bariery.
\end{uwaga}

\section{Funkcjonały Lapunowa i równanie Allena--Cahna}
\begin{twierdzenie}[Dyssypacja energii i stany graniczne]
\label{thm:pde27-allen}
Niech $D,L>0$, $d\ge1$, a
$u_0\in C(\T_L^d)\cap H^1(\T_L^d)$.
Istnieje dokładnie jedno globalne rozwiązanie łagodne
$u$ równania $u_t=D\Delta u+u-u^3$ z danymi $u_0$.
Jest ono klasyczne dla dodatnich czasów.
Dla funkcjonału
\[
 E(u)=\int_{\T_L^d}\left(\frac D2|\nabla u|^2
 +\frac14(u^2-1)^2\right)\,\mathrm d\lambda_d(x)
\]
zachodzi dla $t\ge0$ tożsamość
$E(u(t))+\int_0^t\|u_s(s)\|_2^2\,\mathrm d\lambda_1(s)=E(u_0)$.
Ponadto istnieje ciąg $t_n\to\infty$ i funkcja
$u_\infty\in H^1(\T_L^d)$ takie, że
$u(t_n)\to u_\infty$ silnie w $L^2$ i
$-D\Delta u_\infty=u_\infty-u_\infty^3$
w sensie słabym. Nie twierdzimy bez dodatkowych
założeń o zbieżności całej trajektorii do jednego
stanu stacjonarnego.
\end{twierdzenie}
\begin{proof}
Wybieramy $M\ge\max\{1,\|u_0\|_\infty\}$.
Stałe $-M$ i $M$ są odpowiednio pod- i
nadrozwiązaniem, gdyż $f(-M)=-M+M^3\ge0$
i $f(M)=M-M^3\le0$. Z zasady porównania
otrzymujemy $|u(t,x)|\le M$ i globalność.
Aby uzasadnić tożsamość energii bez dodatkowych
założeń o mieszanych pochodnych, stosujemy
Stwierdzenie~\ref{prop:pde-v44-torus-scaling}.
Oznaczmy $a=D(2\pi/L)^2$, $b=L/(2\pi)$,
$\tau=at$ i $w(\tau,\theta)=u(t,b\theta)$.
Wówczas $w_\tau=\Delta_\theta w-\widetilde V'(w)$,
gdzie $\widetilde V(s)=(s^2-1)^2/(4a)$.
Dane $w(0)$ należą do $C(\T^d)\cap H^1(\T^d)$,
ponieważ liniowe skalowanie zachowuje te przestrzenie.
Stosujemy Twierdzenie~\ref{thm:pde25-energy},
z funkcjonałem
$\widetilde E(w)=\int_{\T^d}
(\frac12|\nabla_\theta w|^2+\widetilde V(w))
\,\mathrm d\lambda_d(\theta)$.
Zmiana zmiennych i reguła łańcuchowa dają
\[
 E(u(t))=ab^d\widetilde E(w(at)),\qquad
 \|u_t(t)\|_2^2=a^2b^d\|w_\tau(at)\|_2^2.
\]
Mnożąc tożsamość energii dla $w$ przez $ab^d$
i podstawiając $\sigma=as$ w całce czasowej,
otrzymujemy dokładnie żądaną tożsamość dla $u$,
również od chwili $0$.
W szczególności
$\|\nabla u(t)\|_2^2\le2E(u_0)/D$ oraz
$\|u(t)\|_2^2\le M^2L^d$. Zatem
$\sup_{t\ge0}\|u(t)\|_{H^1}<\infty$.
Ponieważ energia jest nieujemna,
$\int_0^\infty\|u_t(t)\|_2^2\,\mathrm d\lambda_1(t)
\le E(u_0)$ na mocy zbieżności monotonicznej.
Dla każdego $n\ge1$ wybieramy $t_n\in[n,n+1]$
takie, że
$\|u_t(t_n)\|_2^2\le\int_n^{n+1}\|u_t(s)\|_2^2
\,\mathrm d\lambda_1(s)$.
Jest to możliwe dzięki ciągłości funkcji podcałkowej
dla dodatnich czasów. Prawa strona dąży do zera
z całkowalności na półprostej.
Słaba zwartość ograniczonych ciągów w przestrzeni
Hilberta $H^1$ oraz zwarte zanurzenie
$H^1(\T_L^d)\hookrightarrow L^2(\T_L^d)$
pozwalają po przejściu do podciągu założyć
$u(t_n)\rightharpoonup u_\infty$ w $H^1$ i
$u(t_n)\to u_\infty$ w $L^2$.
Z silnej zbieżności w $L^2$ możemy wybrać dalszy
podciąg zbieżny prawie wszędzie; stąd $|u_\infty|\le M$
prawie wszędzie. Funkcja $f(s)=s-s^3$ jest
lipschitzowska na $[-M,M]$ ze stałą co najwyżej
$1+3M^2$, więc
$\|f(u(t_n))-f(u_\infty)\|_2
\le(1+3M^2)\|u(t_n)-u_\infty\|_2\to0$.
Dla każdej $\varphi\in H^1(\T_L^d)$ równanie
w chwili $t_n>0$ daje
\[
 D\int_{\T_L^d}\nabla u(t_n)\cdot\nabla\varphi\,
                         \mathrm d\lambda_d(x)
 =\int_{\T_L^d}(f(u(t_n))-u_t(t_n))\varphi\,
                         \mathrm d\lambda_d(x).
\]
Lewa strona zbiega dzięki słabej zbieżności
gradientów, prawa dzięki silnej zbieżności w $L^2$
i zanikaniu $\|u_t(t_n)\|_2$.
Otrzymujemy słabą postać równania stacjonarnego
dla $u_\infty$.
\end{proof}

\section{Układy i niestabilność Turinga}
Dla wektorowej funkcji $U=(u_1,u_2)^{\mathsf T}$
rozważamy $U_t=\mathcal D\Delta U+F(U)$,
gdzie $F\in C^1(\R^2;\R^2)$,
$\mathcal D=\operatorname{diag}(D_1,D_2)$, $D_i>0$.
Stan jednorodny $U_*$ spełnia $F(U_*)=0$.
Po podstawieniu $U=U_*+\varepsilon V$
i podzieleniu równania przez $\varepsilon$
składnik reakcji przyjmuje postać
$\varepsilon^{-1}(F(U_*+\varepsilon V)-F(U_*))$.
Z różniczkowalności $F$ jego granicą jest $JV$,
gdzie $J=DF(U_*)$. Linearyzacja ma więc postać
$V_t=\mathcal D\Delta V+JV$.
W odróżnieniu od przypadku skalarnego dyfuzja może
wywołać niestabilność stanu stabilnego dla samego ODE.


\begin{twierdzenie}[Kryterium liniowej niestabilności Turinga]
\label{thm:pde27-turing-general}
Niech $J=\begin{psmallmatrix}a&b\\c&d\end{psmallmatrix}$
będzie rzeczywistą macierzą $2\times2$ oraz
$\mathcal D=\operatorname{diag}(D_1,D_2)$, $D_1,D_2>0$.
Załóżmy $\operatorname{tr}J<0$ i $\det J>0$,
tak że jednorodny układ liniowy $V_t=JV$ jest
asymptotycznie stabilny. Dla $\mu>0$ macierz
$J-\mu\mathcal D$ ma wartość własną o dodatniej
części rzeczywistej wtedy i tylko wtedy, gdy
\begin{equation}\label{eq:pde27-turing-poly}
 D_1D_2\mu^2-(D_2a+D_1d)\mu+\det J<0.
\end{equation}
Takie $\mu>0$ istnieje wtedy i tylko wtedy, gdy
\[
 S:=D_2a+D_1d>0,\qquad S^2>4D_1D_2\det J.
\]
Na torusie lub ograniczonym obszarze Lipschitzowskim
z jednorodnym warunkiem Neumanna występowanie
wykładniczo rosnącego modu jest równoważne temu,
że co najmniej jedna dodatnia wartość własna
$\mu_k$ operatora $-\Delta$ spełnia
\eqref{eq:pde27-turing-poly}; warunki na $S$ i wyróżnik
nie wystarczają bez tego warunku spektralnego.
\end{twierdzenie}
\begin{proof}
Jeżeli $-\Delta\phi=\mu\phi$ i
$V(t,x)=v(t)\phi(x)$, gdzie $v(t)\in\R^2$,
to podstawienie do równania liniowego daje
$v'=(J-\mu\mathcal D)v$. Oznaczamy tę macierz
przez $M_\mu$.
Z hipotez wynika
$\operatorname{tr}M_\mu=\operatorname{tr}J-
\mu(D_1+D_2)<0$. Dla rzeczywistej macierzy $2\times2$
o ujemnym śladzie dodatnią część rzeczywistą
wartości własnej otrzymujemy dokładnie przy
$\det M_\mu<0$: wartości własne są wtedy
rzeczywiste i mają przeciwne znaki.
Jeśli wyznacznik jest dodatni, to rzeczywiste
pierwiastki mają ten sam znak, a ich ujemna suma
wymusza oba znaki ujemne; dla pary zespolonej
sprzężonej część rzeczywista jest połową śladu.
Przy wyznaczniku zerowym pierwiastki to $0$
i $\operatorname{tr}M_\mu<0$. Są różne, więc
mod jest ograniczony, lecz nie asymptotycznie stabilny.
Zatem dodatni wykładnik wzrostu zachodzi dokładnie
dla ujemnego wyznacznika. Rachunek
\[
 \det M_\mu=(a-D_1\mu)(d-D_2\mu)-bc
 =D_1D_2\mu^2-(D_2a+D_1d)\mu+\det J
\]
daje wielomian z~\eqref{eq:pde27-turing-poly}.
Jeżeli $S\le0$, wszystkie składniki trójmianu
są nieujemne dla $\mu>0$, a wyraz wolny jest dodatni.
Jeżeli $S>0$, minimum na dodatniej półosi
jest osiągane w $\mu_0=S/(2D_1D_2)$ i wynosi
$\det J-S^2/(4D_1D_2)$. Jest ujemne dokładnie
przy dodatnim wyróżniku. Wtedy oba pierwiastki
są dodatnie, bo mają dodatnią sumę i iloczyn.

Na torusie $\T_L^d$ funkcje
$\exp(2\pi i k\cdot x/L)$, $k\in\mathbb Z^d$,
tworzą po normalizacji bazę ortonormalną w $L^2$
i odpowiadają wartościom
$\mu_k=(2\pi/L)^2|k|^2$. Dla warunku Neumanna
na ograniczonym obszarze Lipschitzowskim analogiczny
rozkład daje baza funkcji własnych realizacji $-\Delta$
z formą $\int_\Omega\nabla v\cdot\nabla w\,
\mathrm d\lambda_d(x)$ na $H^1(\Omega)$.
Każdy współczynnik ewoluuje niezależnie według
macierzy $M_{\mu_k}$, zatem rosnący wykładniczo
mod istnieje dokładnie przy ujemnym wyznaczniku
dla co najmniej jednej dopuszczalnej wartości $\mu_k$.
Dla dużych $\mu$ zachodzi ponadto
\[
 \tfrac12\frac{\mathrm d}{\mathrm dt}|v|^2
 \le(\|J\|-\mu\min\{D_1,D_2\})|v|^2,
\]
więc wszystkie dostatecznie wysokie mody
zanikają wykładniczo z jednostajnym oszacowaniem.
Pozostało tylko skończenie wiele wartości własnych
(z krotnościami); gdy żaden mod nie rośnie,
ich ewolucje są ograniczone, a jeśli wszystkie
wyznaczniki są dodatnie, zanikają wykładniczo.
Przez tożsamość Parsevala daje to odpowiednio
ograniczoność lub wykładniczy zanik w $L^2$.
\end{proof}

\begin{uwaga}[Przedział niestabilnych częstości przestrzennych]
Jeżeli warunki z poprzedniego twierdzenia są spełnione,
niestałe mody rosną dokładnie dla
$\mu\in(\mu_-,\mu_+)$, gdzie
$\mu_\pm=(S\pm\sqrt{S^2-4D_1D_2\det J})/(2D_1D_2)$.
Na torusie należy więc sprawdzić, czy
$\mu_-<(2\pi/L)^2|k|^2<\mu_+$ dla pewnego $k\ne0$.
Dla jednakowych współczynników dyfuzji $D_1=D_2=D$
mamy $M_\mu=J-\mu DI$, a jego wartości własne
powstają z wartości własnych $J$ przez odjęcie
$\mu D$; destabilizacja dyfuzyjna nie jest możliwa.
W przykładzie poniżej $\mu_-=1/5$ i $\mu_+=1/2$;
wybór $\mu_1=3/10$ rzeczywiście trafia do tego przedziału.
\end{uwaga}


\begin{stwierdzenie}[Kryterium niestabilności konkretnego modu]
\label{prop:pde27-turing}
Niech $\Omega=(0,L)$, na końcach niech obowiązują
jednorodne warunki Neumanna, oraz
\[
 J=\begin{pmatrix}1&-2\\2&-3\end{pmatrix},\qquad
 \mathcal D=\begin{pmatrix}1&0\\0&10\end{pmatrix}.
\]
Macierz $J$ jest stabilna w sensie liniowym dla ODE
$\dot V=JV$: obie jej wartości własne wynoszą $-1$.
Jeżeli $L=\pi/\sqrt{3/10}$, to odpowiedni
jednomodowy układ reakcji--dyfuzji ma dodatnią
rzeczywistą wartość własną. W szczególności
zero jest liniowo niestabilne dla PDE, mimo
stabilności przestrzennie jednorodnej dynamiki.
\end{stwierdzenie}
\begin{proof}
Funkcje $\cos(n\pi x/L)$ są funkcjami własnymi
laplasjanu Neumanna z wartościami $-\mu_n$,
$\mu_n=(n\pi/L)^2$. Na takim modzie ewolucja
jest wyznaczona przez macierz $J-\mu_n\mathcal D$.
Dla $n=0$ mamy
$\operatorname{tr}J=-2$, $\det J=1$, więc
wielomian charakterystyczny jest $(\lambda+1)^2$.
Macierz $J+I$ jest niezerowa i $(J+I)^2=0$,
więc $e^{tJ}=e^{-t}(I+t(J+I))\to0$, także
mimo braku diagonalizowalności.
Dla $n=1$, $\mu_1=3/10$, otrzymujemy
\[
 \operatorname{tr}(J-\mu_1\mathcal D)=-\frac{53}{10},
 \qquad \det(J-\mu_1\mathcal D)
 =1-7\frac3{10}+10\frac9{100}=-\frac15.
\]
Ujemny wyznacznik oznacza jedną dodatnią i jedną
ujemną rzeczywistą wartość własną. Dodatnia wartość własna wynosi
$\lambda_+=(-53+\sqrt{2889})/20>0$.
Jeżeli $v\ne0$ jest odpowiadającym jej rzeczywistym
wektorem własnym, to
$V(t,x)=e^{\lambda_+t}v\cos(\pi x/L)$
jest rzeczywistym rozwiązaniem liniowego PDE
z warunkiem Neumanna, rosnącym wykładniczo.
Ponieważ $\mu_n=(3/10)n^2$, jedynie $n=1$
należy do przedziału $(1/5,1/2)$.
To dowodzi twierdzenia.
\end{proof}

\begin{center}
\begin{tikzpicture}[x=11cm,y=2.1cm,>=Latex]
 \fill[red!10] (.2,-.33) rectangle (.5,0);
 \draw[->] (-.025,0)--(.74,0) node[right] {$\mu$};
 \draw[->] (0,-.36)--(0,1.35) node[above] {$\det M_\mu$};
 \draw[blue!75!black,thick,domain=0:.69,samples=101]
  plot (\x,{10*\x*\x-7*\x+1});
 \draw[densely dashed,gray] (.2,-.30)--(.2,0);
 \draw[densely dashed,gray] (.5,-.30)--(.5,0);
 \fill (.2,0) circle (1.5pt) node[above left] {$\mu_-$};
 \fill (.5,0) circle (1.5pt) node[above right] {$\mu_+$};
 \draw[red!75!black,dashed] (.3,0)--(.3,-.2);
 \fill[red!75!black] (.3,-.2) circle (2pt);
 \node[below,red!75!black] at (.3,-.34) {$\mu_1=3/10$};
 \node[left] at (0,1) {$1$};
\end{tikzpicture}

\small Rysunek: trójmian $\det M_\mu=10\mu^2-7\mu+1$.
Zacieniowany przedział odpowiada ujemnemu wyznacznikowi;
jedyny dodatni mod Neumanna w tym przedziale to $\mu_1$.
\end{center}

\begin{uwaga}[Stabilność liniowa a nieliniowa]
Powyższy rachunek dowodzi niestabilności
\emph{zlinearyzowanego} PDE. Aby wnioskować o
nieliniowej niestabilności układu
$U_t=\mathcal D\Delta U+F(U)$, trzeba dodatkowo
sprecyzować regularność $F$ oraz zastosować
twierdzenie o niestabilności ze spektralnej wartości
własnej o dodatniej części rzeczywistej.
Takiego twierdzenia nie ukrywamy w argumencie liniowym.
\end{uwaga}

\section{Podsumowanie pojęciowe}
Monotoniczność operatorów z rozdziału~\ref{chap:pde26}
prowadziła do istnienia nieliniowych rozwiązań eliptycznych.
W obecnym rozdziale kluczową rolę odgrywają
\emph{porównanie} i \emph{dynamika półgrupy}.
W skalarnej teorii KPP dyfuzja i reakcja prowadzą
do frontu o minimalnej prędkości $2\sqrt{Dr}$;
w układach z różnymi współczynnikami dyfuzji
mogą natomiast powstawać niestabilne mody przestrzenne.
Tych dwóch mechanizmów nie należy utożsamiać.

